% ══════════════════════════════════════════════════════════════════
%  SECTION VII --- RESULTS ANALYSIS AND DISCUSSION
% ══════════════════════════════════════════════════════════════════
\section{Results Analysis and Discussion}
\label{sec:discussion}

\subsection{Hypotheses}
\begin{table}[H]
\centering
\caption{Verdicts on the Phase 2 hypotheses}
\label{tab:verdicts}
\scriptsize
\setlength{\tabcolsep}{3.5pt}
\begin{tabularx}{\textwidth}{>{\raggedright\arraybackslash}p{4.3cm} Y >{\raggedright\arraybackslash}p{2.4cm}}
\toprule
\textbf{Hypothesis} & \textbf{Evidence} & \textbf{Verdict} \\
\midrule
H1: slow explainers make each flow at least 100$\times$ slower & True for 25 of 44 slow settings (up to 1{,}500$\times$); not on the slow Random Forest or with small budgets; fast explainers add only 1.1--49$\times$ & Partly supported \\
H2: explanation quality falls under encryption & CIC-IDS2017: quality \emph{rises} for 7 of 8 pairs; HIKARI: 3 worse, 1 better, 4 unchanged & Rejected on CIC-IDS2017; mixed on HIKARI \\
H3: explanations change across networks, in step with accuracy loss & Below the noise levels in 12 of 16 cases; drift follows the recall drop ($\rho=0.63$) & Supported \\
H4: the trigger keeps overhead $\le$ 10\%, explains $\ge$ 95\% of alerts in time, with no loss of quality & Overhead 0.101--0.106, 100\% of alerts explained within 1\,s, quality equivalent to fresh explanations & Supported within detection capacity \\
\bottomrule
\end{tabularx}
\end{table}

\subsection{Main findings}
\begin{enumerate}[leftmargin=*, itemsep=0pt]
  \item \textbf{Cost depends on the explainer--detector pair, not on ``XAI'' in general.} Fast model-specific explainers cost about as much as detection and are also the most faithful. XGBoost with TreeSHAP, or an EBM, detects and explains a flow for 1.8--4.2 times the cost of detection alone. KernelSHAP and LIME belong only on a background lane.
  \item \textbf{Encryption removes artefacts, not explainability.} On CIC-IDS2017 the $T_3$ view loses at most 1.5 F1 points, gives more faithful explanations, and produces fewer false positives after a network change.
  \item \textbf{The choice of explainer changes the story.} Explainers share only about 2 of their top-5 features, so a security team should fix one explainer and measure its quality on its own traffic.
  \item \textbf{Explanation drift is a free early warning} of a network change.
  \item \textbf{Selective explanation helps most at the edges:} during floods (cache), under overload (pausing) and for missed attacks (uncertainty rule).
  \item \textbf{Padding is a cheap way to evade} a metadata-based detector, so it should be tested.
\end{enumerate}

\subsection{Research gaps}
Gaps G1--G3 were found in Phase~1 and G4 was added in Phase~2. Table~\ref{tab:gapmap} maps each gap to what this phase shows and what remains open.

\begin{table}[!htbp]
\centering
\caption{Research gaps of Phase 2: what Phase 3 shows and what remains open}
\label{tab:gapmap}
\scriptsize
\setlength{\tabcolsep}{3.5pt}
\begin{tabularx}{\textwidth}{>{\raggedright\arraybackslash}p{3.2cm} Y >{\raggedright\arraybackslash}p{3.8cm}}
\toprule
\textbf{Gap} & \textbf{What Phase 3 shows} & \textbf{Still open} \\
\midrule
G1: explanation cost not measured & Time and memory of every explainer on one core; published costs include explainer building (Sections~\ref{sec:rq1}, \ref{sec:compare}) & Edge devices; GPU batching \\
G2: quality under encryption unknown & Quality of nine explainers on three visibility tiers, with paired tests (Section~\ref{sec:rq2}) & Real QUIC captures; packet-sequence models \\
G3: consistency under network change & Drift beyond the noise levels, linked to accuracy loss (Section~\ref{sec:rq3}) & Other shifts (DoH tunnel tools); drift-triggered retraining \\
G4: no budget-aware explanation policy & A trigger tested against seven baselines and fifteen ablation variants (Sections~\ref{sec:rq4}, \ref{sec:ablation}) & Adaptive budgets; a study with real analysts \\
\bottomrule
\end{tabularx}
\end{table}

\subsection{Strengths and limitations}
\textbf{Strengths.} Cost, quality and drift are measured on the same explanations of the same flows; the data are leakage-controlled with time-based splits; the encryption assumption is tested on real packets; the explainers are checked against a known answer; every comparison uses paired statistical tests with effect sizes; and every result can be regenerated from the code.

\textbf{Limitations.}
\begin{itemize}[leftmargin=*, itemsep=0pt]
  \item The QUIC view is emulated by removing columns from flow datasets, not measured on real QUIC traffic, and the only truly encrypted dataset (HIKARI) has weak detectors.
  \item All timings come from one laptop, with 200 flows per setting and one seed.
  \item Removal-based quality metrics depend on the value that replaces a removed feature; two different replacement values and a known-answer test (Section~\ref{sec:abl:sanity}) gave the same explainer ranking.
  \item Many statistical tests were run, so all of them are Holm-corrected.
\end{itemize}
